\pagestyle{myheadings} \setcounter{page}{1} \setcounter{footnote}{0}

\section{~设置嵌套运行} \label{app:nest}
\newcounters
\vssub
\subsection{~使用 {\file ww3\_shel}}
\vssub

使用单网格波浪模式程序 {\file ww3\_shel} 运行嵌套模式，其机制原则上很
简单。一个大尺度模式生成包含边界数据的文件，例如 {\file nest1.ww3}。
随后将该文件重命名为 {\file nest.ww3}，并放入运行嵌套（小尺度）模式的
目录中。小尺度模式随后会自动处理该文件，并在需要且可用时更新边界条件。
一致地设置嵌套则更复杂。这里给出一种简单的分步方法。下一小节还介绍了
另一种可能：使用 {\code ww3\_bound} 从谱输出组装 {\file nest.ww3} 文件。

\newcounter{nestlist}

\begin{list}{\arabic{nestlist})}{\usecounter{nestlist}
             \rightmargin 5mm \labelsep 3mm
             \labelwidth 5mm \leftmargin 10mm}

\item 第一步是完整设置大尺度模式，但暂不为嵌套模式生成边界数据。包括
      适当的风场、图形输出等。测试该模式，直到确认其运行正常。

\item 设置小尺度模式，暂时忽略边界条件。需要考虑的是，理想情况下边界
      条件应与大尺度模式中的网格线重合，以尽量减小边界数据文件大小。
      以与大尺度模式相同的方式设置该模式，并充分测试。

\item 当小尺度模式按上述方式设置满意后，需要定义边界条件。进入小尺度
      模式的 {\file ww3\_grid.inp} 文件，按第~\ref{sub:ww3grid} 节
      文档中的说明标记所有预期输入边界。确保在 {\file switch} 文件中
      选择了模式开关 {\F !/O1}，必要时重新编译。运行 {\code ww3\_grid}
      并保存屏幕输出。该程序的输出现在包含所有被标记为输入边界点的列表。
      还要确保小尺度模式的 {\file mod\_def.ww3} 存储副本（如有）已正确
      更新。

\item 下一步是在大尺度模式中将上述列表里的所有输入边界点包含为输出边界
      点。保留该列表，转到大尺度模式的 {\file ww3\_grid.inp} 文件。按
      第~\ref{sub:ww3grid} 节文档中的说明，将上述列表中的所有点加入为
      输出边界点。确保所有数据（且没有其他数据）写入单个文件，并用适当
      的输入文件运行 {\code ww3\_grid}。现在应得到一个输出边界点列表，
      它应与上述输入边界点列表一致。注意，点在列表中出现的顺序并不重要。
      同样要确保大尺度模式的 {\file mod\_def.ww3} 存储副本（如有）已正确
      更新。

\item 如果两个点列表之间存在差异，在前两个步骤之间迭代，直到列表一致。

\item 下一步是开始从大尺度模式生成边界数据。这需要在大尺度模式中激活
      嵌套输出。输出已经在上述步骤中设置，并包含在大尺度模式的模式定义
      文件（{\file mod\_def.ww3}）中。现在需要在大尺度模式实际运行的输入
      文件 {\file ww3\_shel.inp} 中设置起始时间、时间增量和结束时间来激活
      它。如果添加的是第二个或后续嵌套，则无需执行此步骤。大尺度模式现在
      会生成包含边界数据的文件。如果这是第一个包含的嵌套，输出文件将为
      {\file nest1.ww3}。该文件需要保存，以供小尺度模式使用。

\item 为了在小尺度模式中使用嵌套数据，上述边界数据文件需要重命名为
      {\file nest.ww3}，并放入运行小尺度模式 {\file ww3\_shel} 的目录。
      如果小尺度模式已在其定义文件 {\file mod\_def.ww3} 中正确定义输入
      边界点，它会自动处理 {\file nest.ww3} 文件，并在边界数据可用时进行
      更新。此时建议再做两个测试。

\begin{itemize}

\item 首次在存在 {\file nest.ww3} 文件的情况下运行小尺度模式时，密切关注
      {\code ww3\_shel} 的输出，确认 (i) 程序报告已处理 {\file nest.ww3}
      文件并判定其正常，且 (ii) 没有关于边界数据不兼容或缺失的额外警告。
      还应检查日志文件 {\file log.ww3}，确认边界数据在预期时间更新。

\item 当所有数据看起来都已处理后，较为直观且稳妥的做法是运行一次小尺度
      模式，其中在 {\file ww3\_shel.inp} 中关闭风场，并且不提供重启动文件
      {\file restart.ww3}。在这种模式运行中，波浪能量只能从边界进入计算
      域。这是一个很好的测试，用于确认边界数据按预期从大尺度模式传递到
      小尺度模式。

\end{itemize}

\end{list}

\noindent
可以用同样方式添加更多嵌套模式。从小尺度模式添加第二层嵌套也同样进行。
模式目前设置为每次模式运行最多生成 9 个边界数据文件。连续（“望远镜式”）
嵌套的数量没有限制。

\vssub
\subsection{~使用 {\file ww3\_bound} 和/或非结构网格}
\vssub

在某些情况下，运行大尺度模式时很难或不可能预先知道小尺度模式强迫点的
位置。如果希望使用来自在线数据库或第三方数据库的边界条件来运行海岸放大
网格，就会出现这种情况。

在这种情况下，可以使用 {\code ww3\_bound} 从谱输出生成 {\file nest.ww3}
文件。由于边界上点的间距不规则，这对非结构网格也特别方便。
{\code ww3\_bound} 接收一个谱文件列表，这些文件应具有相同的谱网格，并
生成一个 {\file nest.ww3}。插值系数由最近可用谱的位置和小尺度模式中活动
边界点的位置确定。

\pb
\subsection{~使用 {\file ww3\_multi}}
\vssub

与使用波浪模式驱动程序 {\file ww3\_shel} 相比，在波浪模式驱动程序
{\file ww3\_multi} 中执行双向嵌套大为简化，因为所需的所有数据传输均由
多网格波浪模式例程在内部完成。镶嵌模式系统通过反复执行以下步骤来设置。

\begin{list}{\arabic{nestlist})}{\usecounter{nestlist}
             \rightmargin 5mm \labelsep 3mm
             \labelwidth 5mm \leftmargin 10mm}

\item 使用 {\file ww3\_grid} 工具设置一个网格。定义该网格、其活动边界点
      以及时间步长等所有其他模式信息，但{\it 不要}尝试为其他网格生成输出
      嵌套数据。这将由 {\file ww3\_multi} 中的多网格波浪模式例程自动评估。
      注意，最低等级网格可选择使用活动边界数据，这些数据可以从文件读取，
      也可以在计算期间保持常值。较高等级网格需要活动边界点，才能在镶嵌
      方法中有效，

\item 将该网格作为额外网格加入 {\file ww3\_multi.inp} 输入文件，并指定
      合适的等级号。运行 {\file ww3\_multi} 将识别网格与请求的边界数据点
      之间可通过迭代解决的差异，以及网格之间的其他差异。手工去除这类差异
      可能很繁琐。\citet{tol:MMAB07a, tol:OMOD08a} 的网格生成包会自动检查
      这类差异，因此推荐用于本版本 \ws\ 的网格生成。

\end{list}

\noindent
注意，定义输入数据场的网格也可以用类似方式加入。还要注意，建议在海洋
输入场（流、水位和冰）中使用陆海掩膜，以确保海岸点处具有现实的输入值。

通常先开发较低等级网格，不过任意等级的网格都可以在任意时间加入。

%\bpage \pagestyle{empty}
